%% kapa: exam=1 round=36 year=2025 period=1 form=A subject=civil_law subject_ko=민법 seq=20 number=20
%% source: Q-Net 감정평가사 기출문제 (cst00309) artlSeq=5231525 file=raw/1008/5231525/2025년 제36회 감정평가사 1차시험 1교시 문제지 원본.pdf
%% conversion: kordoc 4.24.51 -> tools/gen1.py (auto); review=proofread-sonnet; flags=none
\begin{question}{20}
불공정한 법률행위에 관한 설명으로 옳지 않은 것은? (다툼이 있으면 판례에 따름)
\begin{choices}
  \item 매매계약이 불공정한 법률행위에 해당하는지는 계약체결 당시를 기준으로 판단해야 한다.
  \item 궁박은 심리적 원인에 기인할 수도 있다.
  \item 무경험은 특정영역에서의 경험부족을 뜻한다.
  \item 대리행위의 경우에 궁박은 본인을 기준으로 판단한다.
  \item 현저하게 공정을 잃었는지는 거래상의 객관적 가치에 따라 판단해야 한다.
\end{choices}
\end{question}
\answer{3}
